\documentclass[10pt,a4paper]{amsart}
\usepackage[margin=19mm]{geometry}
\usepackage{amsmath,amssymb,mathtools}
\usepackage[scr=rsfs,cal=euler]{mathalfa}
\usepackage{xfrac}
\usepackage[unicode,hidelinks,bookmarks=false]{hyperref}
\newcommand{\ie}{i.e.\ }
\newcommand{\cf}{cf.\ }
\newcommand{\resp}{respectively\ }
\newcommand{\Subsection}{\subsection}
\providecommand{\nobreakdash}{\nobreak\hbox{-}}
\setlength{\emergencystretch}{2em}
\multlinegap0pt
\usepackage{amssymb}
\usepackage{braket}
\usepackage{bm}
\usepackage{dsfont}
\usepackage[normalem]{ulem}
\usepackage{cancel}
\usepackage{xcolor}
\usepackage{mathtools}
\newtheorem{lemma}{Lemma}
\title{Critical point search and linear response theory for computing electronic excitation energies of molecular systems. \\ $\mbox{Part I: General framework, application to Hartree-Fock and DFT}$}
\author{Laura Grazioli, Yukuan Hu, Eric Canc\`es}
\date{Source: arXiv:2506.16420v2 (CC BY 4.0)}
\begin{document}
\maketitle

\noindent\emph{Source and licence.} Critical point search and linear response theory for computing electronic excitation energies of molecular systems. \\ $\mbox{Part I: General framework, application to Hartree-Fock and DFT}$. Version arXiv:2506.16420v2 (CC BY 4.0), distributed by arXiv under the Creative Commons Attribution 4.0 International licence, http://creativecommons.org/licenses/by/4.0/, as recorded in the arXiv licence metadata for this exact version and shown on the abstract page https://arxiv.org/abs/2506.16420v2. Full licence notice: \texttt{licenses/arxiv-2506.16420-CC-BY-4.0.txt}.\par\smallskip

\noindent\textit{Editorial scope.} Counted unit: the article's lemma lemma (source0.tex lines 1155--1171). The article's complete original proof of it is delivered with it (source0.tex lines 1172--1189, one \texttt{proof} environment of 18 lines). No mathematics is written, repaired or re-derived: the statement and the proof are the article's own lines, in the article's own notation, and no retained line is substituted or re-flowed.  No further source-local context is needed: every cross-reference of every retained line resolves inside the two delivered spans.  Nothing was omitted from any retained span, so the package carries no omission record.  No retained line contains a citation, so the wrapper carries no bibliography and no bibliography entry.  No retained line contains a figure environment, an image-inclusion command or drawing code, so no image is delivered and the package declares no asset dependency.  Editorial formatting only, in addition: an AMS \texttt{amsart} preamble that restates the article's own declarations and macros used by the retained lines, namely the environments \texttt{lemma} on separate counters, together with the standard packages those lines need.  The retained environments print in this document as Lemma 1 in order of appearance, whereas the article prints the same items with the numbering of its own full document.  The complete native source of the article as distributed in the arXiv e-print for this exact version is bundled unchanged as \texttt{sources/180155/source0.tex}.
\begin{lemma}
    Consider a family $(A(\eta))_{\eta \ge 0}$ of real symmetric $(4n) \times (4n)$ matrices of the form
    $$
A(\eta) = \left( \begin{array}{cc} A_+(\eta) & 0 \\ 0 & A_-(\eta) \end{array} \right), \quad \mbox{with} \quad A_\pm(\eta)=\Omega_0 + \eta B_\pm + \mathcal O(\eta^2), 
    $$
    and 
    $$
    \Omega_0 = {\rm diag}(\omega_1^0,\cdots \omega_n^0,\omega_1^0,\cdots,\omega_n^0) \quad \mbox{with} \quad 0 < \omega_1^0 < \cdots < \omega_n^0 .
    $$
    If the two eigenvalues $\lambda_{i,\pm}$ of the real symmetric matrix
    $$\frac12\begin{pmatrix}
        [B_++B_-]_{ii} & [B_++B_-]_{i(n+i)}\\
        [B_++B_-]_{(n+i)i} & [B_++B_-]_{(n+i)(n+i)}
    \end{pmatrix}$$
    are distinct, then for $\eta \ge 0$ small enough, the symplectic eigenvalues of $A(\eta)$ can be expanded as
    $$\omega_{i,\pm}(\eta)=\omega_i^0+\eta\lambda_{i,\pm}+\mathcal O(\eta^2),\quad i=1,\ldots,n.$$
\end{lemma}

\begin{proof}
    For $\eta \ge 0$ small enough, the symplectic eigenvalues of $A(\eta)$ are the eigenvalues of the symmetric matrix
    $$
    \Omega(\eta):=\left( A_-(\eta)^{1/2} A_+(\eta) A_-(\eta)^{1/2} \right)^{1/2}.
    $$
    Denoting by $\omega_{n+i}^0:=\omega_i^0$ for $1 \le i \le n$, we have
    \begin{align*}
    & A_-(\eta)^{1/2} = \Omega_0^{1/2} + \eta C + \mathcal O(\eta^2) \quad \mbox{with} \quad [C]_{ij}=\frac{[B_-]_{ij}}{(\omega_i^0)^{1/2}+(\omega_j^0)^{1/2}}, \\
    & A_-(\eta)^{1/2} A_+(\eta) A_-(\eta)^{1/2} = \Omega_0^2 + \eta D + \mathcal O(\eta^2), \quad \mbox{with} \quad D = \Omega_0^{3/2}C + \Omega_0^{1/2} B_+ \Omega_0^{1/2} + C \Omega_0^{3/2}, \\
    & \Omega(\eta) = \Omega_0 + \eta \Omega_1 + \mathcal O(\eta^2), \quad \mbox{with} \quad [\Omega_1]_{ij} = \frac{[D]_{ij}}{\omega_i^0+\omega_j^0}.
    \end{align*} 
    By simple algebraic manipulation, we obtain
    \begin{align*}
      [\Omega_1]_{ij} &=  \frac{1}{\omega_i^0+\omega_j^0} \left( \frac{(\omega_i^0)^{3/2} + (\omega_j^0)^{3/2}}{(\omega_i^0)^{1/2}+(\omega_j^0)^{1/2}} [B_-]_{ij} + (\omega_i^0)^{1/2} (\omega_j^0)^{1/2} [B_+]_{ij} \right) \\
      & = [B_-]_{ij} +  \frac{(\omega_i^0)^{1/2}(\omega_j^0)^{1/2}}{\omega_i^0+\omega_j^0}  [B_+-B_-]_{ij}.
    \end{align*}
   The proof is completed after applying degenerate first-order perturbation theory.
\end{proof}

\end{document}
